\section{Notation}\label{section-notation}

Let ${\mathcal{E}}_{ij}$, $i,j=1\lcd n$, be the standard generators of the Lie algebra $\gl_n$, with the commutation relations
\begin{gather*}
[{\mathcal{E}}_{ij},{\mathcal{E}}_{kl}]=\delta_{jk}{\mathcal{E}}_{il}-\delta_{il}{\mathcal{E}}_{kj} ,
\end{gather*}
where $\delta_{jk}$ is the Kronecker symbol. We shall also use the root notation ${\mathcal{H}}_\alpha$, ${\mathcal{E}}_{\alpha}$, ${\mathcal{E}}_{-\alpha}$, \dots\ for elements of~$\gl_n$.


Let ${\mathcal{E}}_{ij}^{(1)}$ and ${\mathcal{E}}_{ij}^{(2)}$, $i,j=1\lcd n$, be the standard generators of the two copies of the Lie algebra $\gl_n$ in $\g:=\gl_n\oplus\gl_n$,
\begin{gather*}
[{\mathcal{E}}_{ij}^{(a)},{\mathcal{E}}_{kl}^{(b)}]=\delta_{ab}\big(\delta_{jk}{\mathcal{E}}_{il}^{(a)}-\delta_{il}{\mathcal{E}}_{kj}^{(a)}\big) .\end{gather*}
Set
\begin{gather*}
 e_{ij}:={\mathcal{E}}_{ij}^{(1)}+{\mathcal{E}}_{ij}^{(2)} ,\qquad E_{ij}:={\mathcal{E}}_{ij}^{(1)}-{\mathcal{E}}_{ij}^{(2)}
 .
 \end{gather*}
The elements $e_{ij}$ span the diagonally embedded Lie algebra $\k\simeq\gl_n$, while $E_{ij}$ form an adjoint $\k$-module $\mathfrak{p}$. The
Lie algebra $\k$ and the space $\mathfrak{p}$ constitute a symmetric pair, that is, $[\k,\k]\subset \k$, $[\k,\p]\subset \p$, and $[\p,\p]\subset \k$:
\begin{gather*}
 [e_{ij},e_{kl}] = \delta_{jk}e_{il}-\delta_{il}e_{kj}  ,\qquad
  [e_{ij},E_{kl}] = \delta_{jk}E_{il}-\delta_{il}E_{kj}  ,\qquad
   [E_{ij},E_{kl}] = \delta_{jk}e_{il}-\delta_{il}e_{kj} .
\end{gather*}
In the sequel, $h_a$ means the element $e_{aa}$ of the Cartan subalgebra $\h$ of the subalgebra $\k\in\gl_n\oplus\gl_n$ and $h_{ab}$ the
element $e_{aa}-e_{bb}$.


Let $\{\ve_a\}$ be the basis of $\h^*$ dual to the basis $\{ h_a\}$ of $\h$, $\ve_a(h_b)=\delta_{ab}$. We shall use as well the root notation $h_\alpha$,
$e_{\alpha}$, $e_{-\alpha}$ for elements of $\k$, and $H_\alpha$, $E_{\alpha}$, $E_{-\alpha}$ for elements of $\p$.


The Lie subalgebra $\nplus$ in the triangular decomposition \rf{intro1} is spanned by the root vectors~$e_{ij}$ with $i<j$ and the Lie subalgebra $\nminus$ by the
root vectors $e_{ij}$ with $i>j$. Let $\bb_+$ and $\bb_-$
be the corresponding Borel subalgebras, $\bb_+=\h\oplus\nplus$ and $\bb_-=\h\oplus\nminus$. Denote by $\Delta_+$ and $\Delta_-$ the sets of positive and
negative roots in the root system $\Delta=\Delta_+\cup \Delta_-$ of $\k$: $\Delta_+$ consists of roots $\ve_i-\ve_j$ with $i<j$ and $\Delta_-$ consists of roots
$\ve_i-\ve_j$ with $i>j$.
Let $\Q$ be the root lattice, $\Q:=\{\gamma\in\h^*\,|\,\gamma=\sum_{\alpha\in\Delta_+, n_\a\in\ZZ}n_\a \a\}$. It contains the
positive cone $\Q_+$,
\begin{gather*}
\Q_+:=\bigg\{\gamma\in\h^*\,|\,\gamma= \sum_{\alpha\in\Delta_+, n_\a\in\ZZ, \; n_\a\geq 0}n_\a \a\bigg\} .%\label{cqp}
\end{gather*}
For $\lambda,\mu\in\h^*$, the notation
\begin{gather}
\lambda>\mu\label{paor}
\end{gather}
means that the dif\/ference $\lambda-\mu$ belongs to~$\Q_+$, $\lambda-\mu\in\Q_+$. This is a partial order in~$\h^*$.


We f\/ix the following action of the cover of the symmetric group $\S_n$ (the Weyl group of the
diagonal $\k$) on the Lie algebra $\gl_n\oplus\gl_n$ by automorphisms
\begin{gather*}
\sy_i(x):=\mathrm{Ad}_{\exp(e_{i,i+1})}\mathrm{Ad}_{\exp(-e_{i+1,i})}\mathrm{Ad}_{\exp(e_{i,i+1})}(x) ,
\end{gather*}
so that
\begin{gather*}\sy_i(e_{kl})=(-1)^{\delta_{ik}+\delta_{il}}e_{\sigma_i(k)\sigma_i(l)} ,
\qquad\sy_i(E_{kl})=(-1)^{\delta_{ik}+\delta_{il}}E_{\sigma_i(k)\sigma_i(l)}.
\end{gather*}
Here $\sigma_i=(i,i+1)$ is an elementary transposition in the symmetric group. We extend naturally the above action of the cover of $\S_n$ to the action by
automorphisms on the associative algebra $\Ar\equiv\Ar_n:=\U(\gl_n)\otimes \U(\gl_n)$. The restriction of this action to $\h$ coincides
with the natural action $\si(h_{k})=h_{\si(k)}$, $\si\in\S_n$,  of the Weyl group on the Cartan subalgebra.


Besides, we use the shifted action of $\S_n$ on the polynomial algebra $\U(\h)$ (and its localizations) by automorphisms; the shifted action is def\/ined by
\begin{gather}
\label{not1}
\si\circ h_k:= h_{\si(k)}+k-\si(k) ,\qquad k=1,\dots ,n;\qquad \si\in\S_n .
\end{gather}
It becomes the usual action for the variables
\begin{gather}
 \th_k:=h_k-k ,\qquad \th_{ij}:=\th_i-\th_j; \label{hrond}
\end{gather}
by \rf{not1} for any $\si\in\S_n$ we have
\begin{gather*}
\si\circ\th_k=\th_{\si(k)} ,\qquad \si\circ\th_{ij}=\th_{\si(i)\si(j)} .
\end{gather*}

It will be sometimes convenient to denote the commutator $[a,b]$ of two elements $a$ and $b$ of an associative algebra by
\begin{gather}
 \hat{a}(b):=[a,b] .\label{dehaco}
\end{gather}


\section[Reduction algebra $\Z_n$]{Reduction algebra $\boldsymbol{\Z_n}$}\label{section2}

In this section we recall the def\/inition of the reduction algebras, in particular the diagonal reduction algebras of the $\gl$ type.
We introduce the order for which
the ordering relations for the algebra $\Z_n$ will be discussed.
The formulas for the Zhelobenko automorphisms for the algebra $\Z_n$ are given; some basic facts about the standard involution, anti-involution
and central elements for the algebra $\Z_n$ are presented at the end of the section.

{\bf 1.} Let $\Uh$ and $\Ab$ be the rings of fractions of the algebras $\U(\h)$ and $\Ar$ with respect to the multiplicative set, generated by elements
\begin{gather*}
 h_{ij}+l ,\qquad l\in\ZZ  , \quad 1\leq i<j\leq n  .%\label{musel}
 \end{gather*}


Def\/ine $\Z_n$ to be the double coset space of $\Ab$ by its left ideal $\Ib:=\Ab\nplus$, generated by elements of $\nplus$, and the right ideal
$\Jb:=\nminus\Ab$, generated by elements of $\nminus$, $\Z_n:=\Ab/(\Ib+\Jb)$.


The space $\Z_n$ is an associative algebra with respect to the multiplication map
\begin{gather}
\label{not5a}
a\mult b:=aP b  .
\end{gather}
Here $P$ is the extremal projector \cite{AST} for the diagonal $\gl_n$.
It is an element of a certain extension of the algebra $\U(\gl_n)$ satisfying the relations $e_{ij}P=Pe_{ji}=0$ for all $i$ and
$j$ such that $1\leq i<j\leq n$.


The algebra  $\Z_n$ is a particular example of a reduction algebra; in our context, $\Z_n$ is
def\/ined by the coproduct (the diagonal inclusion) $\U(\gl_n)\to\Ar$.

{\bf 2.}  The main structure theorems for the reduction algebras are given in \cite[Section~2]{KO3}.


In the sequel we choose a weight linear basis $\{p_K\}$ of $\p$ ($\p$ is the $\k$-invariant complement to~$\k$ in~$\g$,
$\g =\k +\p$) and equip it with a total order $\prec$. The total order~$\prec$ will be compatible with the partial order $<$ on
$\h^*$, see~\rf{paor}, in the sense that $\mu_K<\mu_L\ \Rightarrow \ p_K\prec p_L$. We shall sometimes write
$I\prec J$ instead of $p_I\prec p_J$. For an arbitrary element $a\in\Ab$ let $\widetilde{a}$ be its image in the reduction
algebra; in particular, $\pp_K$ is the image in the reduction algebra of the basic vector $p_K\in\p$.

{\bf 3.} In our situation we choose the set of vectors $E_{ij}$, $i,j=1,\dots ,n$, as a basis of the space~$\p$. The weight of $E_{ij}$ is $\ve_i-\ve_j$.
The compatibility of a total order $\prec$ with the partial order~$<$ on~$\h^*$ means the condition
\begin{gather*}%\label{not4a}
E_{ij}\prec E_{kl}\qquad \text{if}\quad i-j>k-l .
\end{gather*}
The order in each subset $\{E_{ij}|i-j=a\}$ with a f\/ixed $a$ can be chosen arbitrarily. For instance, we can set
\begin{gather}\label{not4}
E_{ij}\prec E_{kl}\qquad \text{if}\quad i-j>k-l\quad\text{or}\quad i-j=k-l\quad\text{and}\quad i>k .
\end{gather}

Denote the images of the elements $E_{ij}$ in $\Z_n$ by $\s_{ij}$. We use also the notation $t_i$ for the elements $\s_{ii}$ and $t_{ij}:=t_i-t_j$ for the elements
$\s_{ii}-\s_{jj}$. The order \rf{not4} induces as well the order on the generators $\s_{ij}$ of the algebra $\Z_n$:
\begin{gather*} %\label{not4b}
\s_{ij}\prec\s_{kl}\ \Leftrightarrow\ E_{ij}\prec E_{kl}  .
\end{gather*}
The  statement (d) in the paper \cite[Section~2]{KO3} implies an existence of structure constants $\mathrm{B}_{(ab),(cd),(ij),(kl)} \in \Uh$ and $\mathrm{D}_{(ab),(cd)}\in\Uh$
such that for any $a,b,c,d=1,\ldots,n$ we have
\begin{gather}\label{not6}
\s_{ab}\mult\s_{cd}=\sum_{i,j,k,l:\s_{ij}\preceq\s_{kl}}\mathrm{B}_{(ab),(cd),(ij),(kl)}\s_{ij}\mult\s_{kl}+\mathrm{D}_{(ab),(cd)}  .
\end{gather}
In particular, the algebra $\Z_n$ (in general, the reduction algebra related to a symmetric pair $(\k,\p)$, $\g :=\k +\p$) is $\ZZ_2$-graded; the degree of $\s_{ab}$ is~1 and the degree of any element from $\Uh$ is~0.


The relations \rf{not6} together with the weight conditions
\begin{gather*}
 [h,\s_{ab}]=(\ve_a-\ve_b)(h) \s_{ab}%\label{wede}
\end{gather*}
are the def\/ining relations for the algebra $\Z_n$.


Note that the denominators of the structure constants $\mathrm{B}_{(ab),(cd),(ij),(kl)}$ and
 $\mathrm{D}_{(ab),(cd)}$ are pro\-ducts
of linear factors of the form $\th_{ij}+\ell$, $i<j$, where $\ell\geq -1$ is an integer, see~\cite{KO3}.

{\bf 4.} The algebra $\Z_n$ can be equipped with the action of Zhelobenko automorphisms \cite{KO}.  Denote by $\q_{i}$  the Zhelobenko
automorphism $\q_i:\Z_n\to\Z_n$  corresponding to the transposition $\si_i\in\S_n$. It is def\/ined as follows~\cite{KO}. First we def\/ine a map
$\q_i:\Ar\to \Ab/\Ib$ by
\begin{gather} \label{not7}
\q_i(x):=\sum_{k\geq 0}\frac{(-1)^k}{k!} \hat{e}_{i,i+1}^k(\sy_i(x)) e_{i+1,i}^k   \prod_{a=1}^k(h_{i,i+1}-a+1)^{-1}
\quad \mod\Ib  .
\end{gather}
Here $\hat{e}_{i,i+1}$ stands for the adjoint action of the element $e_{i,i+1}$, see~\rf{dehaco}.
The operator $\q_i$ has the property
\begin{gather}
\label{not2}
\q_i(hx)=(\si_i\circ h)\q_i(x)
\end{gather}
for any $x\in\Ar$ and  $h\in\h$; $\si\circ h$ is def\/ined in~\rf{not1}. With the help of~(\ref{not2}), the map $\q_i$ can be extended
to the map (denoted by the same symbol) $\q_i:\Ab\to \Ab/\Jb$ by
setting $\q_i(a(h)x)=(\si_i\circ a(h))\q_i(x)$ for any $x\in\Ar$ and
$a(h)\in\Uh$. One can further prove that $\q_i(\Ib)=0$ and $\q_i(\Jb)\subset (\Jb+\Ib)/\Ib$,
so that $\q_i$ can be viewed as a linear operator $\q_i:\Z_n\to\Z_n$. Due to~\cite{KO},
this is an algebra automorphism, satisfying~\rf{not2}.


The operators $\q_i$ satisfy the braid group relations~\cite{Zh}:
\begin{gather*}
\q_i\q_{i+1}\q_i=\q_{i+1}\q_{i}\q_{i+1}  ,\\
\q_i\q_j=\q_j\q_i  ,\quad |i-j|>1  ,
\end{gather*}
and the inversion relation \cite{KO}:
\begin{gather}
\label{invr}
\q_i^2(x)=\frac{1}{h_{i,i+1}+1}  \sy_i^2(x)  (h_{i,i+1}+1)  ,\qquad x\in\Z_n  .
\end{gather}
In particular, $\q_i^2(x)=x$ if $x$ is of zero weight.

{\bf 5.} The Chevalley anti-involution $\epsilon$ in $\U(\gl_n\oplus\gl_n)$,
$\epsilon(e_{ij}):=e_{ji}$, $\epsilon(E_{ij}):=E_{ji}$, induces the anti-involution $\epsilon$ in
the algebra $\Z_n$:
\begin{gather}
\epsilon(\s_{ij})=\s_{ji}  ,\qquad \epsilon(h_k)=h_k  .\label{anep}
\end{gather}
Besides, the outer automorphism of the Dynkin diagram of $\gl_n$ induces the involutive automorphism
$\omega$ of $\Z_n$,
\begin{gather}
\label{not2a}
\omega(\s_{ij}) =(-1)^{i+j+1}\s_{j'i'}  ,\qquad \omega (h_k)=-h_{k'}  ,
\end{gather}
where $i'=n+1-i$. The operations $\epsilon$ and $\omega$ commute, $\epsilon\omega =\omega\epsilon$.


Central elements of the subalgebra $\U(\gl_n)\otimes 1\subset \Ar$, generated by $n$ Casimir
operators of degrees $1\lcd n$,  as well as central elements of the subalgebra
$1\otimes \U(\gl_n)\subset \Ar$ project to central elements of the algebra $\Z_n$. In particular,
central elements of degree $1$ project to central elements
\begin{gather}
I^{(n,h)}:=h_1+\dots +h_n\label{clih}
\end{gather}
and
\begin{gather}
\label{clit} I^{(n,t)}:=t_1+\dots +t_n
\end{gather}
of the algebra $\Z_n$. The dif\/ference of central elements of degree two projects to the central element
\begin{gather}
\label{drclit}
\sum_{i=1}^n(h_i-2i)t_i
\end{gather}
of the algebra $\Z_n$. The images of other Casimir operators are more complicated.


\section{Main results}

This section contains the principal results of the paper. We f\/irst give
preliminary information on the new basis in which the def\/ining relations for the
algebra $\Z_n$ can be written down in an economical fashion. The braid group action on the new generators is then
explicitly given in Subsection~\ref{brgac}. The complete set of the def\/ining relations for the algebra $\Z_n$ is written down in Subsection~\ref{section3.3}. The regime for which both
the set of the derived def\/ining relations and the set of the def\/ining ordering relation have a controllable ``limiting behavior'' is introduced in Subsection~\ref{seclimit}. Subsection~\ref{subsection_sl_n} deals with the diagonal reduction algebra for $\sl_n$; the quadratic Casimir operator for
${\mathrm{DR}}(\sl_n)$ as well as for the diagonal reduction algebra for an arbitrary semi-simple
Lie algebra~$\k$ is given there. Subsection~\ref{subsection3.4} is devoted to the stabilization and cut phenomena with respect to the embedding of the Lie algebra $\gl_n\oplus\gl_m$ into
the Lie algebra~$\gl_{n+m}$; the theorem about the behavior of the centers of the diagonal reduction algebra under the cutting is proved there.

\subsection{New variables}

We shall use the following elements of $\Uh$:
\begin{gather*}  \tphi_{ij}:=\frac{\th_{ij}}{\th_{ij}-1}  ,\quad \tpsi_{ij}:=\frac{\th_{ij}-1}{\th_{ij}} ,
\quad \ttphi_{ij}:=\frac{\th_{ij}-1}{\th_{ij}-2} ,\quad \ttpsi_{ij}:=\frac{\th_{ij}-2}{\th_{ij}-1} ,\quad \Cprime_{ij}:=\frac{\th_{ij}-3}{\th_{ij}-2},
\end{gather*}
the variables $\th_{ij}$ are def\/ined in~\rf{hrond}. Note that $\tphi_{ij}\tpsi_{ij}=\ttphi_{ij}\ttpsi_{ij}=1$.


Def\/ine  elements $\tt_1,\ldots, \tt_n\in\Z_n$ by
\begin{gather*}
\label{3.1}
\tt_1:=t_1 ,\quad\tt_2:=\q_{1}(t_1) ,\quad\tt_3:=\q_{2}\q_{1}(t_1) ,
\quad\ldots, \quad \tt_n:=\q_{n-1}\cdots \q_{2}\q_{1}(t_1)  .
\end{gather*}
Using \rf{not7} we f\/ind the relations
 \begin{gather}
 \begin{split}
& \q_i(t_i)  =-\frac{1}{\th_{i,i+1}-1}t_i+\frac{\th_{i,i+1}}{\th_{i,i+1}-1}t_{i+1}  ,
\qquad \q_i(t_{i+1})= \frac{\th_{i,i+1}}{\th_{i,i+1}-1}t_i-\frac{1}{\th_{i,i+1}-1}t_{i+1} ,\\ %\nonumber\\
& \q_i(t_k) =  t_k ,\qquad k\not=i,i+1  ,
 \end{split}\label{acwot}
 \end{gather}
which can be used to convert the def\/inition \rf{3.1} into  a linear over the ring $\Uh$ change of variables:
\begin{gather}
\begin{split}
& \ds{\tt_l} =\ds{t_l\prod_{j=1}^{l-1}\tphi_{jl}
-\sum_{k=1}^{l-1}t_k\frac{1}{\th_{kl}-1}\prod_{j=1}^{k-1}\tphi_{jl}}  , \\
& \ds{t_l} =\ds{\tt_l\prod_{j=1}^{l-1}\tpsi_{jl} +\sum_{k=1}^{l-1}\tt_k\frac{1}{\th_{kl}}
\prod_{\substack{j=1 \\ j\neq k}}^{l-1}\tpsi_{jk}} .
\end{split} \label{3.1a}
\end{gather}
For example,
\begin{gather*}
\tt_2 =-\frac{1}{\th_{12}-1}t_1+\frac{\th_{12}}{\th_{12}-1}t_2 ,\qquad
t_2=\frac{1}{\th_{12}}\tt_1+\frac{\th_{12}-1}{\th_{12}}\tt_2 ,\\
\tt_3 =-\frac{1}{\th_{13}-1}t_1-\frac{\th_{13}}{(\th_{13}-1)(\th_{23}-1)}t_2+
\frac{\th_{13}\th_{23}}{(\th_{13}-1)(\th_{23}-1)}t_3  ,\\
t_3 =  \frac{\th_{12}+1}{\th_{12}\th_{13}}\tt_1+\frac{\th_{12}-1}{\th_{12}\th_{23}}\tt_2+
\frac{(\th_{13}-1)(\th_{23}-1)}{\th_{13}\th_{23}}\tt_3  .
\end{gather*}


In terms of the new variables $\tt$'s, the linear in $t$ central element (\ref{clit}) reads
\begin{gather*}
\sum t_i=\sum\tt_i\prod_{a:a\neq i}\frac{\th_{ia}+1}{\th_{ia}} .
\end{gather*}

\subsection{Braid group action}
\label{brgac}

Since
$\q_{i}^2(x)=x$ for any element $x$ of zero weight, the braid group acts as its symmetric group quotient on the space of weight 0 elements.
It follows from \rf{3.1} and $\q_{i}(t_1)=t_1$ for all $i>1$ that
\begin{gather}\label{3.2}
\q_{\sigma}(\tt_i)=\tt_{\sigma(i)}
\end{gather}
for any $\sigma\in \S_n$.


The action of the Zhelobenko automorphisms, see Section \ref{section2},  on the generators $\s_{kl}$  looks as follows:
\begin{alignat}{3}
& \q_i(\s_{ik})=-\s_{i+1,k}\tphi_{i,i+1} , \qquad  &&
\q_i(\s_{ki})=-\s_{k,i+1} ,\quad k\not=i,i+1 ,& \nonumber\\
&\q_i(\s_{i+1,k})=\s_{i,k} ,\qquad && \q_i(\s_{k,i+1})=\s_{k,i}\tphi_{i,i+1} ,\quad k\not=i,i+1 ,& \label{3.5}\\
& \q_i(\s_{i,i+1})=-\s_{i+1,i}\tphi_{i,i+1}\ttphi_{i,i+1} ,\qquad && \q_i(\s_{i+1,i})=-\s_{i,i+1} , & \nonumber\\
&\q_{i}(\s_{j,k})=\s_{j,k} ,\quad j,k\neq i,i+1 .\qquad &&&\nonumber
\end{alignat}

Denote $i'=n+1-i$, as before. The braid group action \rf{3.5} is compatible with the anti-involution $\epsilon$
and the involution $\omega$ (note that $\omega (\th_{ij})=\th_{j'i'}$),  see~\rf{anep} and~\rf{not2a}, in the following sense:
\begin{gather}
\epsilon  \q_i =\q_i^{-1}\epsilon  ,\\
 \omega  \q_i =\q_{i'-1}\omega .\label{qeps}
\end{gather}


Let $w_0$ be the longest element of the Weyl group of $\gl_n$, the symmetric group $\S_n$. Similarly to the squares of
the transformations corresponding to the simple roots, see \rf{invr}, the action of~$\q_{w_0}^2$ is the conjugation by a
certain element of $\Uh$.

\begin{lem}\label{lqw02}
We have
\begin{gather}\label{qw02}
\q_{w_0}^2(x)=S^{-1}xS  ,
\end{gather}
where
\begin{gather}
 S=\prod_{i,j:i<j}\th_{ij}  .\label{conjq02}
 \end{gather}
\end{lem}

The proof shows that the formula \rf{qw02} works for an arbitrary reductive Lie algebra, with $S=\prod_{\alpha\in \Delta_+}\th_{\alpha}$.

\begin{prop}\label{paqw0}
The action of $\q_{w_0}$ on generators reads
\begin{gather}
 \label{3.6}
 \q_{w_0}(\s_{ij})=(-1)^{i+j} \s_{i'j'}\prod_{a:a<i'}\tphi_{ai'}
\prod_{b:b>j'}\tphi_{j'b} ,\\
\q_{w_0}(\tt_i)=\tt_{i'} .\label{3.6a}
\end{gather}
\end{prop}

The proofs of Lemma \ref{lqw02} and Proposition \ref{paqw0} are in Section \ref{sectionproofs}.

\subsection{Def\/ining relations}\label{section3.3}

To save  space we omit in this section the symbol $\mult$ for the multiplication in the algebra~$\Z_n$. It should not lead to any confusion since no other multiplication is used in this section.


Each relation which we will derive will be of a certain weight, equal to a sum of two roots. From general considerations the upper estimate
for the number of terms in a quadratic relation of weight $\lambda=\alpha+\beta$ is the number $|\lambda|$ of quadratic combinations $\s_{\alpha'}\s_{\beta'}$
with $\alpha'+\beta'=\lambda$.
There are several types, excluding the trivial one, $\lambda=2(\ve_i-\ve_j)$, $|\lambda|=1$:
\begin{enumerate}\itemsep=0pt
\item $\lambda=\pm(2\ve_i-\ve_j-\ve_k)$, where $i,j$ and $k$ are pairwise distinct. Then $|\lambda|=2$.
\item $\lambda=\ve_i-\ve_j+\ve_k-\ve_l$ with pairwise distinct $i,j,k$ and $l$. Then $|\lambda|=4$.
\item $\lambda=\ve_i-\ve_j$, $i\neq j$. For $\s_{\alpha'}\s_{\beta'}$, there are $2(n-2)$ possibilities (subtype 3a) with $\alpha'=\ve_i-\ve_k$,
$\beta'=\ve_k-\ve_j$ or $\alpha'=\ve_k-\ve_j$, $\beta'=\ve_i-\ve_k$ with $k\neq i,j$ and $2n$ possibilities (subtype 3b) with $\alpha'=0$,
$\beta'=\ve_i-\ve_j$ or $\alpha'=\ve_i-\ve_j$, $\beta'=0$. Thus $|\lambda|=4(n-1)$.
\item $\lambda=0$. There are $n^2$ possibilities (subtype 4a) with $\alpha'=0$, $\beta'=0$ and $n(n-1)$ possibilities (subtype 4b) with
$\alpha'=\ve_i-\ve_j$,  $\beta'=\ve_j-\ve_i$, $i\neq j$. Here $|\lambda|=n(2n-1)$.
\end{enumerate}

Below we write down relations for each type (and subtype) separately. The relations of the types 1 and 2 have a simple form in terms of the original
generators $\s_{ij}$. To write the relations  of the types 3 and 4, it is convenient to renormalize the generators $\s_{ij}$
with $i\not=j$. Namely, we set
\begin{gather}
\label{not8}\hs_{ij}=\s_{ij}\prod_{k=1}^{i-1}\tphi_{ki}  .
\end{gather}

In terms of the generators $\hs_{ij}$, the formulas \rf{3.5} for the action of the automorphisms $\q_i$  translate as follows:
\begin{alignat*}{3}
& \q_{i}(\hs_{ik})=-\hs_{i+1,k} ,\qquad &&  \q_{i}(\hs_{i+1,k})=\hs_{i,k}\tphi_{i+1,i} ,\quad k\not=i,i+1 , & \\
%\label{not9}
&\q_{i}(\hs_{ki})=-\hs_{k,i+1} ,\qquad && \q_{i}(\hs_{k,i+1})=\hs_{k,i}\tphi_{i,i+1}=\tpsi_{i+1,i}\hs_{k,i} ,\quad k\not=i,i+1,&
\\
&\q_{i}(\hs_{i,i+1})=-\tpsi_{i+1,i}\hs_{i+1,i} , \qquad &&\q_{i}(\hs_{i+1,i})=-\hs_{i,i+1}\tphi_{i+1,i} ,& \\
 &\q_{i}(\hs_{j,k})=\hs_{j,k} , \quad j,k\neq i,i+1 .\qquad &&&
 \end{alignat*}

{\bf 1.} The relations of  the type 1 are:
\begin{gather}\label{relation1}
\s_{ij}\s_{ik}=\s_{ik}\s_{ij}\tphi_{kj} ,\qquad
\s_{ji}\s_{ki}=\s_{ki}\s_{ji}\tpsi_{kj} ,\qquad \text{for}\quad j<k ,\  i\not=j,k .
\end{gather}

{\bf 2.} Denote
\begin{gather*}
D_{ijkl}:=\left( {\ds\frac{1}{\th_{ik}}}-{\ds\frac{1}{\th_{jl}}}\right).
\end{gather*}
Then, for any four pairwise dif\/ferent indices $i$, $j$, $k$ and $l$, we have the following relations of the type 2:
\begin{gather}
\begin{split}
& [ \s_{ij},\s_{kl} ]=\s_{kj}\s_{il}D_{ijkl} , \qquad i<k ,\ j<l  , \\ %\nonumber\\
& \s_{ij}\s_{kl}-\s_{kl}\s_{ij}\tphi_{jl}'\tphi_{lj}'=\s_{kj}\s_{il}
D_{ijkl} ,\qquad i<k ,\ j>l.
\end{split} \label{relation2}
\end{gather}

{\bf 3a.} Let $i\neq k\neq l\neq i$. Denote
\begin{gather*}
\mathring{E}_{ikl}:=-\left((\tt_i-\tt_k) \frac{\th_{il}+1}{\th_{ik}\th_{il}}+(\tt_k-\tt_l) \frac{\th_{il}-1}{\th_{kl}\th_{il}}\right)\hs_{il}+\sum_{a:a\neq i,k,l}
\hs_{al}\hs_{ia} \frac{\ttphi_{ai}}{\th_{ka}+1} .
\end{gather*}
With this notation the f\/irst group of the relations of the type 3 is:
\begin{gather}\notag\hs_{ik}\hs_{kl}\tpsi_{ik} -\hs_{kl}\hs_{ik}\ttphi_{ki}
=\mathring{E}_{ikl} , \qquad i<k<l ,\\   \notag
\hs_{ik}\hs_{kl}\tpsi_{ik}\tpsi_{lk}\ttphi_{lk} -\hs_{kl}\hs_{ik}\ttphi_{ki}
=\mathring{E}_{ikl}  , \qquad i<l<k ,\\
\label{relation3}
\hs_{ik}\hs_{kl}\tphi_{ki} -\hs_{kl}\hs_{ik}\ttphi_{ki}=\mathring{E}_{ikl} , \qquad  k<i<l ,\\
\notag \hs_{ik}\hs_{kl}\tphi_{ki}\tphi_{li}\ttpsi_{li} -\hs_{kl}\hs_{ik}\ttphi_{ki}
=\mathring{E}_{ikl}  ,\qquad  k<l<i ,\\
\notag \hs_{ik}\hs_{kl}\tpsi_{ik}\tpsi_{lk}\ttphi_{lk}\tphi_{li}\ttpsi_{li}-\hs_{kl}\hs_{ik}\ttphi_{ki}
 =\mathring{E}_{ikl}  , \qquad l<i<k , \\  \notag
\hs_{ik}\hs_{kl}\tphi_{ki}\tpsi_{lk}\ttphi_{lk}\tphi_{li}\ttpsi_{li}
-\hs_{kl}\hs_{ik}\ttphi_{ki}=\mathring{E}_{ikl} , \qquad l<k<i .
\end{gather}

The relations \rf{relation3} can be written in a more compact way with the help of both systems, $\s_{ij}$ and $\hs_{ij}$, of generators. Let now
\begin{gather*}
E_{ikl}:=-\left((\tt_i-\tt_k) \frac{\th_{il}+1}{\th_{ik}\th_{il}}+(\tt_k-\tt_l) \frac{\th_{il}-1}{\th_{kl}\th_{il}}\right)
\s_{il}+\sum_{a:a\neq i,k,l}
\hs_{al}\s_{ia} \frac{\ttphi_{ai}}{\th_{ka}+1} .
\end{gather*}
Then
\begin{gather}
\begin{split}
& \s_{ik}\hs_{kl}\tpsi_{ik} -\hs_{kl}\s_{ik}\ttphi_{ki}=E_{ikl} ,\qquad k<l,\\ %\nonumber\\
& \s_{ik}\hs_{kl}\tpsi_{ik}\tpsi_{lk}\ttphi_{lk} -\hs_{kl}\s_{ik}\ttphi_{ki}=E_{ikl} ,\qquad l<k .
\end{split}
\label{shfo3a}
\end{gather}
Moreover, after an extra redef\/inition: $\!\mathring{\hspace{.1cm}z}_{kl}\!\!\!\!\!\!\mathring{\phantom{z}}\ =\hs_{kl}\ttphi_{lk}$ for $k>l$, the left hand side of the
second line in \rf{shfo3a} becomes, up to a common factor, the same as the left hand side of the f\/irst line, namely, it reads $(\s_{ik}
\mathring{\hspace{.1cm}z}_{kl}\!\!\!\!\!\!\mathring{\phantom{z}}\ \ \tpsi_{ik} -\mathring{\hspace{.1cm}z}_{kl}\!\!\!\!\!\!\mathring{\phantom{z}}\ \ \s_{ik}\ttphi_{ki})
\tpsi_{lk}$.

{\bf 3b.} Let $i\neq j\neq k\neq i$. The second group of relations of the type 3 reads:
\begin{gather}\notag\hs_{ij}\tt_{i}=  \tt_i\hs_{ij}\Cprime_{ji}-\tt_j\hs_{ij}  \frac{1}{\th_{ij}+2}-
\sum_{a:a\not=i,j}\hs_{aj}\hs_{ia}  \frac{1}{\th_{ia}+2}  ,\\
\label{3.12}\hs_{ij}\tt_{j}= -\tt_i\hs_{ij}  \frac{\Cprime_{ji}}{\th_{ij}-1}+
\tt_j\hs_{ij}\tphi_{ij}\tpsi_{ji}\ttphi_{ji}+
\sum_{a:a\not=i,j}\hs_{aj}\hs_{ia}\tphi_{ij}\tpsi_{ji}{\ttphi_{ai}}{\th_{ja}+1}  ,\\  \notag
\hs_{ij}\tt_{k}= \tt_i\hs_{ij}  \frac{(\th_{ij}+3)\ttphi_{ji}}{(\th_{ik}^2-1)(\th_{jk}-1)}+
\tt_j\hs_{ij}  \frac{(\th_{ij}+1)\ttphi_{ji}}{(\th_{ik}-1)(\th_{jk}-1)^2}+
\tt_k\hs_{ij}\tphi_{ik}\tphi_{ki}\tphi_{jk}\ttpsi_{jk}\\
  \notag
  \phantom{\hs_{ij}\tt_{k}=}{}
- \hs_{kj}\hs_{ik}  \frac{(\th_{ij}+1)\ttphi_{ki}}{(\th_{ik}-1)(\th_{jk}-1)}
-\sum_{a:a\not=i,j,k}\hs_{aj}\hs_{ia}  \frac{\th_{ij}+1}{(\th_{ik}-1)(\th_{jk}-1)}
%\cdot
\frac{\ttphi_{ai}}{\th_{ka}+1}  .
\end{gather}

{\bf 4a.} The relations of the weight zero (the type~4) are also divided into 2 groups.
This is the f\/irst group of the relations:
\begin{gather}
 \label{relation4a}[\tt_i,\tt_j]=0.
 \end{gather}
As follows from the proof, the relations \rf{relation4a} hold for the diagonal reduction algebra for an arbitrary reductive
Lie algebra: the images of the generators, corresponding to the Cartan subalgebra, commute.

{\bf 4b.} Finally, the second group of the relations of the type 4 is
\begin{gather}
 \label{relation4}
 [\hs_{ij},\hs_{ji}]=\th_{ij}-\frac{1}{\th_{ij}}(\tt_i-\tt_j)^2+
\sum_{a:a\not=i,j}\left(\frac{1}{\th_{ja}+1}\hs_{ai}\hs_{ia}-\frac{1}{\th_{ia}+1}\hs_{aj}\hs_{ja}\right) ,
\end{gather}
where $i\neq j$.

{\bf Main statement.} Denote by $\mathfrak{R}$  the system \rf{relation1},  \rf{relation2}, \rf{relation3},
\rf{3.12},  \rf{relation4a} and \rf{relation4} of the relations.

\begin{theorem} \label{mthe}
The relations $\mathfrak{R}$ are the defining relations for the weight generators $\s_{ij}$ and
$t_i$ of the algebra $\Z_n$. In particular, the set  \eqref{not6} of ordering relations follows over $\Uh$ from $($and is equivalent to$)$  $\mathfrak{R}$.
\end{theorem}

The derivation of the system $\mathfrak{R}$ of the relations is given in Section~\ref{sectionproofs}. The validity in $\Z_n$ of relations from the
set~$\mathfrak{R}$, together with the results
from~\cite{KO3}, completes the proof of Theorem~\ref{mthe} (Section~\ref{proofmain}).

